\documentclass{utarticle}
\usepackage{hyperref,graphics,xcolor}
\usepackage{fontspec,unicode-math}
\setmainfont[Ligatures={Rare,Historic,TeX}]{LibertinusSerif-Regular.otf}
\setsansfont[Ligatures={Rare,Historic,TeX}]{LibertinusSans-Regular.otf}
\setmathfont{LibertinusMath-Regular.otf}

\begin{document}
\section{foo τέχνη}
This is a pen, not a street. ちょっとcheckしちゃった．
とりあえず\textsf{SANS serif}ゴシックでも．

\cjkxcode`\γ=15\cjkxcode`\ゔ=0\cjkxcode`\ö=6
\kchardef\hoge=`\ゔ
\achardef\fuga=`\ö

\section{barあいうえお}
Gödel ゔ\hoge\fuga\kchar`\2 λ\achar`\λ\\
$\displaystyle \int_0^α a\gamma γ\sqrt{\frac{uあ}{φ(x∧y)}}\,dx$ vs\ 
\cjkxcode`\γ=7
$\displaystyle \int_0^α a\gamma γ\sqrt{\frac{uあ}{φ(x∧y)}}\,dx$.
\parbox<y>{10zw}{%
えーっと，「\textcolor{blue}{ある程度は動いて}いる」ような気がします？
\kansuji1234567890
}

\cjkxcode`\朝=4
ほげら明朝
ほげらゴシック
ほげら丸ゴシック

】【【：：

\parbox<y>{10zw}{%
\XeTeXcharclass`\：=3
\jfont in jis\g=goth10\g
】【【：：

\XeTeXinterchartokenstate=1
】【【：：

\XeTeXinterchartoks 3 3 {\relax}
】【【：：
}

\end{document}
